\documentclass[11pt,a4paper]{article}
\usepackage[T1]{fontenc}
\usepackage[utf8]{inputenc}
\usepackage{lmodern,microtype}
\usepackage[margin=25mm,headheight=14pt]{geometry}
\usepackage{amsmath,amssymb,amsthm,mathtools}
\usepackage{booktabs,longtable,tabularx,array}
\usepackage{listings,xcolor,xurl,enumitem,fancyhdr}
\usepackage[hidelinks,pdfusetitle]{hyperref}
\newcolumntype{L}[1]{>{\raggedright\arraybackslash}p{#1}}
\definecolor{codeback}{gray}{0.97}
\lstset{basicstyle=\ttfamily\footnotesize,backgroundcolor=\color{codeback},
 breaklines=true,columns=fullflexible,keepspaces=true,showstringspaces=false,
 frame=single,rulecolor=\color{black!18},framerule=0.3pt,xleftmargin=3pt,xrightmargin=3pt,
 aboveskip=9pt,belowskip=9pt}
\setlength{\parskip}{5pt}
\setlength{\parindent}{0pt}
\setcounter{tocdepth}{2}
\newcommand{\code}[1]{\texttt{\detokenize{#1}}}
\newcommand{\OO}{\mathcal O}
\newcommand{\pin}{651f2029d0b2cd4da9e4dfdf1f4275a124d23b99}
\newcommand{\file}[1]{\nolinkurl{#1}}
\newtheorem{proposition}{Proposition}[section]
\newtheorem{lemma}[proposition]{Lemma}
\pagestyle{fancy}
\fancyhf{}\fancyhead[L]{\small Asymptotic: incremental source audit}
\fancyhead[R]{\small 10 September 2026}\fancyfoot[C]{\thepage}
\title{Asymptotic: Regularity Contracts and Structural Work}
\author{An incremental technical audit for Vladimir Reshetnikov}
\date{10 September 2026}
\begin{document}
\begin{titlepage}
\vspace*{10mm}
{\Large\scshape Incremental source audit\par}
\vspace{9mm}
{\Huge\bfseries Asymptotic\par}
\vspace{4mm}
{\LARGE Regularity contracts\par and structural work\par}
\vspace{10mm}
{\large VladimirReshetnikov/Asymptotic\par}
\vspace{3mm}
\textbf{Audited commit}\par
{\small\ttfamily\pin\par}
\vspace{3mm}
10 September 2026
\vspace{10mm}

\textbf{Assessment.} The inspected code contains a derivative-contract propagation
problem in the ordinary absolute-value observable path. A real, smooth input
remainder can produce an absolute value with infinitely many cusps accumulating
at the expansion point, while the result still advertises a classical derivative
bound. Two recent repair paths also warrant targeted performance changes: repeated
affine recognition during recursive interval evaluation, and repeated membership
queries when finding discarded truncation blocks.

\begin{tabularx}{\textwidth}{@{}lXl@{}}
\toprule ID & New contribution & Evidence class\\\midrule
N01 & Classical differentiability is not preserved by the pure-remainder \code{Abs} path. & Source + proof\\
N02 & A linear affine recognizer is repeatedly invoked on nested nonlinear subtrees. & Source + exact counts\\
N03 & Prefix truncation is processed as repeated general membership searches. & Source + cost model\\
\bottomrule
\end{tabularx}

\vspace{4mm}
\textbf{What ran.} Fourteen independent mathematical and structural tests and six
patch-emitter fixture tests passed. They are not Wolfram package tests. Native
Wolfram and Mathics execution did not succeed in this environment; the archive
contains unexecuted native observations and conservative candidate edits for
local validation.

\vfill
\textbf{Scope discipline.} This article does not reissue the existing review
backlog. It uses the consolidated four-wave register and intake records as its
exclusion map, and directly compares the relevant portions of review 14. The ZIP
contains this article, editable \LaTeX, models, recorded independent results,
and reproduction tools. It is not an upstream package release.
\end{titlepage}

\begingroup
\small
\setlength{\parskip}{0pt}
\tableofcontents
\endgroup
\newpage

\section{Baseline, exclusions, and evidence}
\subsection{The reviewed version}
The audit is pinned to commit \code{651f2029d0b2cd4da9e4dfdf1f4275a124d23b99},
whose timestamp is 10 September 2026, 13:42:02 UTC. The relevant library is
\code{AsymptoticAnalysis}; its canonical implementation is under
\file{src/Kernel/}, while the root \file{AsymptoticAnalysis.wl} is a generated
standalone. Findings refer to this immutable snapshot, not to whatever
\code{main} contains when the article is read.\cite{repo,layout}

The immediately preceding commit, \code{e9142f7f5e3db241d5ffb8951c72606ed0c89aa9},
adds the rational-affine interval path discussed in N02. Its purpose is a valid
repair of the earlier C20 problem. N02 concerns the cost of that newly installed
path on \emph{nonlinear} syntax; it does not allege that the original affine
correctness/availability repair is still missing.\cite{affinecommit,cert}

The source was read through pinned GitHub file requests. There was no successful
complete local checkout. Source coverage was deliberately concentrated on the
ordinary jet calculus, source transformations, exact interval evaluation,
several special-function adapters, and selected Mathics compatibility modules.
Appendix~\ref{app:coverage} distinguishes substantial file reads from selected
windows. An inspected path without a retained new finding is not certified
bug-free.

\subsection{How duplication was avoided}
The repository contains 36 imported review packages. Its register describes 231
attributed entries before consolidation of overlaps. Those are not 231 distinct
open defects. The primary exclusion sources here are
\file{CODE_REVIEW_STATUS.md}, the wave-3 and wave-4 intake records, and the
imported-review index. Review 14's affine and truncation sections were additionally
read directly because they are the antecedents of N02 and N03.\cite{reviewindex,status,wave3,wave4,review14}

This is not a claim to have reread every page of all 36 articles. It is an
explicitly documented crosscheck against the maintained consolidation, with
closer original-source comparison where a candidate overlaps earlier work.
The article does not repeat the existing recommendations about native-tail
admission, generic derivative hypotheses, generic resource budgets, backward
precision planning, or quantitative-bound transport as new findings.

\begin{longtable}{@{}L{12mm}L{44mm}L{\dimexpr\textwidth-56mm-4\tabcolsep\relax}@{}}
\toprule New ID & Nearest existing entries & Distinct contribution retained here\\\midrule\endhead
N01 & C13; D02; X04 & A specific \code{Abs} operation carries an already legitimate derivative contract into a result for which that classical derivative does not exist on any punctured neighborhood. The witness is real; no missing initial derivative hypothesis is used.\\
N02 & C20; review 14 N02 & The installed repair repeatedly searches overlapping nonlinear subtrees. The new claim is a quadratic recognizer-call count for a stable Horner syntax family, not failure to preserve affine cancellation.\\
N03 & C22; review 14 N05 & The installed bound-transport routine uses repeated membership queries even when the retained rows are an identical prefix. The new proposal changes this selection algorithm, not the bound-transport theorem.\\
\bottomrule
\end{longtable}

\subsection{Evidence labels and execution limits}
Three different kinds of statement occur below. A \emph{source observation}
identifies a concrete definition and control-flow path in the pinned code. A
\emph{mathematical or structural result} is proved independently of any kernel
execution. A \emph{native reproduction candidate} is executable Wolfram code
whose actual outcome still needs to be recorded on the pinned package.

The independent Python tests exercise the counterexample identities, rational
interval models, structural operation counts, prefix/fallback selection, and
patch-emitter mechanics. They do not emulate Wolfram evaluation, package loading,
option resolution, or Mathics. In particular, a passing independent test does not
mean the public package call has been run. The native service returned connection
errors, and neither a local Wolfram kernel nor an installable Mathics runtime was
available. No native timing, patched native test, or full upstream suite result is
claimed.

N01 is therefore a source-derived regularity-contract defect with a rigorous
admissible witness, rather than an observed incorrect numerical coefficient.
N02's repeated recursive recognizer work is explicit in the implementation.
N03's comparison counts describe a conventional linear-membership model; they
are not a statement about undocumented optimizations inside a particular
Wolfram \code{MemberQ} implementation.

\section{The invariant that matters: a jet is more than its finite part}
The ordinary representation is usefully read as
\begin{equation}\label{eq:representation}
 F(y)=o(y)+p(y)\bigl(J(w(y))+R(w(y))\bigr),\qquad w\to0^+,
\end{equation}
where the finite jet has the form
\[
 J(w)=\sum_{\alpha<P}w^\alpha C_\alpha(\log w),\qquad
 R(w)=\OO\bigl(w^P(1+|\log w|)^D\bigr).
\]
The ordinary \code{SeriesRepresentation} stores an offset, a prefactor, a
positive scale coordinate, a logarithmic variable, and a jet triple. Its
\code{RemainderDerivativeOrder} is a separate contract, used by
\code{seriesDerivative}; a magnitude estimate alone does not satisfy that
guard.\cite{operations,kernel}

For the examples in this article, $o=0$, $p=1$, and $w=y$, so no coordinate
Jacobian or carrier regularity is involved. An order-$k$ derivative contract
means that the actual remainder has the relevant classical derivatives on an
eventual punctured interval and that, for $0\le j\le k$, they obey the matching
bounds
\begin{equation}\label{eq:derivative-contract}
 R^{(j)}(y)=\OO\bigl(y^{P-j}(1+|\log y|)^D\bigr).
\end{equation}
This interpretation is exactly what permits a routine to differentiate a
stored remainder by lowering its power and its available derivative order.
The implementation explicitly rejects a derivative request when the available
order is too small.\cite{operations}

There are two independent closure questions. Does an operation preserve the
\emph{magnitude} estimate? Does it preserve \emph{classical regularity and the
matching derivative estimates}? Absolute value illustrates why the answers can
differ. On real values it is Lipschitz:
\[
 \bigl||a|-|b|\bigr|\le |a-b|.
\]
Thus $|\OO(y^P)|=\OO(y^P)$ is a sound magnitude transport. It does not follow that
absolute value preserves differentiability through zeros of its argument.
No issue about complex absolute value or the interpretation of a fractional
power is needed here.\cite{wlabs}

\section{N01: absolute value preserves a derivative contract it has not proved}
\subsection{The implementation path}
The ordinary helper \code{fwdAbs} starts with the following branch:
\begin{lstlisting}
fwdAbs[j : {T_, P_, D_}, ell_, ass_] := Module[...,
  If[T === {}, Return[j, Module]];
  ...]
\end{lstlisting}
For a pure remainder it returns exactly the same magnitude jet. That is the
right magnitude estimate. For a nonempty jet, the rest of the routine attempts
to establish the eventual sign from the leading logarithmic coefficient.
\cite{kernel}

The expression interpreter \code{seriesJetApply} delegates an \code{Abs}
application to that helper. The public \code{SeriesObservable} then constructs
the result using
\begin{lstlisting}
j = seriesJetApply[body, x, d["Jet"], d, h, limit];
result = seriesMake[Join[d, <|"Jet" -> j|>],
  {"Observable", {s}, e, x}, h];
\end{lstlisting}
The replacement changes the jet but leaves the input's derivative-order field
in \code{d}. The ordinary result constructor retains that field for an inexact
jet. Consequently, the pure-remainder branch can preserve a positive derivative
order without establishing smoothness of the absolute value. The relevant
source locations are \file{SeriesOperations.wl}, the expression interpreter and
lines 354--374, followed by the derivative guard at lines 536--558.
\cite{operations}

\subsection{A legitimate input family, not a forged object}
Consider, for $0<y<1/4$,
\begin{equation}\label{eq:H}
 R(y)=y^3\sin(\log y),\qquad H(y)=y+R(y).
\end{equation}
Every derivative of $R$ exists at every positive $y$. The elementary recurrence
\begin{equation}\label{eq:recurrence}
 \frac{d}{dy}\left[y^m(a\sin\log y+b\cos\log y)\right]
 =y^{m-1}\left[(ma-b)\sin\log y+(a+mb)\cos\log y\right]
\end{equation}
shows that $R^{(j)}(y)=\OO(y^{3-j})$ for every fixed $j$. In particular,
\begin{align}
 R'(y)&=y^2(3\sin\log y+\cos\log y),\label{eq:g}\\
 R''(y)&=5y(\sin\log y+\cos\log y).
\end{align}
A declaration of two matching derivatives for this input remainder is true.

\begin{proposition}\label{prop:inverse-admissible}
The function $H$ in \eqref{eq:H} is the local inverse of a real smooth function
$F$ satisfying $F(x)=x+\OO(x^3)$ on a positive one-sided neighborhood of zero.
\end{proposition}
\begin{proof}
Since $|R'(y)|\le\sqrt{10}\,y^2<4/16$, we have $H'(y)>3/4$ on the stated
interval. Also
\[
 \frac{15}{16}\le\frac{H(y)}y\le\frac{17}{16}.
\]
Thus $H$ is positive and strictly increasing, and has a smooth inverse $F$ on
its positive image. With $x=H(y)$,
\[
 |F(x)-x|=|y-H(y)|\le y^3
 \le \left(\frac{16}{15}\right)^3 x^3.
\]
This proves the required forward remainder estimate and local branch
admissibility.
\end{proof}

The point of this construction is that the input can be an ordinary inverse
result with a declared omitted forward remainder. There is no need to manufacture
an association, exploit an unvalidated saved object, assert false smoothness,
or allow a complex unknown error. The family supplies an actual smooth inverse
realizing the magnitude and derivative contracts used below. The argument uses
the documented magnitude-plus-matching-derivative interpretation of those
contracts, not an extra requirement of analyticity at the endpoint itself.

\subsection{Where the classical derivative disappears}
Set $g=R'$. Write
\[
 g(y)=\sqrt{10}\,y^2\sin(\log y+\theta),\qquad
 \theta=\arctan(1/3).
\]
For every positive integer $n$,
\begin{equation}\label{eq:zeros}
 y_n=\exp(-\theta-n\pi)
\end{equation}
is a zero of $g$, and $y_n\to0^+$. These zeros are simple. Substitution into
$g'(y)=5y(\sin\log y+\cos\log y)$ gives
\begin{equation}\label{eq:slopes}
 |g'(y_n)|=\sqrt{10}\,y_n>0.
\end{equation}

\begin{proposition}\label{prop:cusp}
The magnitude estimate $|g(y)|=\OO(y^2)$ is valid, but $|g|$ has no classical
derivative at any $y_n$. Therefore it is not differentiable on any complete
punctured neighborhood $(0,\delta)$ of zero.
\end{proposition}
\begin{proof}
At a simple zero $y_n$,
$g(y_n+h)=g'(y_n)h+o(h)$. Hence the left and right derivatives of $|g|$ are
$-|g'(y_n)|$ and $+|g'(y_n)|$, respectively. They differ by
$2\sqrt{10}\,y_n$. Every punctured neighborhood contains such a point.
\end{proof}

This is not a failure of a derivative \emph{bound} at points where a derivative
exists. Away from the zeros, the first derivative of $|g|$ is indeed
$\OO(y)$. It is a failure of the classical derivative's existence on the
claimed eventual domain. Quietly interpreting the answer as an almost-everywhere
or weak derivative would change the operation's meaning.

\subsection{The public reproduction candidate}
The following sequence uses the public APIs. The actual native results are not
pre-populated in the archive.
\begin{lstlisting}
s = AsymptoticInverse[x, {x, 0}, {y, 2},
  "InputRemainder" -> {3, 0}];
id = AsymptoticExpansion[y, {y, 0, 4}, "Backend" -> "Package"];
r = SeriesAdd[s, SeriesMultiply[id, -1]];
g = SeriesDifferentiate[r, 1, "RemainderDerivativeOrder" -> 2];
a = SeriesObservable[g, Abs[z], z];
SeriesDifferentiate[a]
\end{lstlisting}
The intended progression, derived from the inspected source, is:
\begin{center}
\begin{tabular}{@{}lll@{}}
\toprule Stage & Magnitude jet & Derivative order\\\midrule
\code{r} & $0+\OO(y^3)$ & Initially undeclared\\
\code{g} & $0+\OO(y^2)$ & $1$ after the truthful order-$2$ declaration\\
\code{a} & $0+\OO(y^2)$ & Incorrectly retained as $1$\\
Final differentiation & $0+\OO(y)$ & Accepted on that retained premise\\
\bottomrule
\end{tabular}
\end{center}
For the member $H$ of the permitted inverse family, \code{a} represents $|g|$
from Proposition~\ref{prop:cusp}. The final classical differentiation must not be
justified by the input contract alone. A corrected implementation should either
refuse the last call, retain an explicitly weaker regularity status, or require
additional information proving an eventual smooth branch.

The proof does not claim that the package executes the unspecified forward
remainder as a numerical function. It reasons about the family denoted by the
remainder contract. The existing numerical checker explicitly distinguishes its
stored explicit equation from an unspecified input-remainder family; this audit
preserves that distinction.\cite{numerical}

\subsection{A small conservative repair}
A minimal correctness repair is to stop propagating a positive derivative order
through an inexact ordinary observable containing \code{Abs}, unless the
operation has separately proved a smooth local branch. The supplied patch emitter
implements the conservative first step immediately after \code{seriesJetApply}:
\begin{lstlisting}
If[! FreeQ[body, _Abs] && j[[2]] =!= Infinity,
  d = Join[d, <|"RemainderDerivativeOrder" -> 0|>]];
\end{lstlisting}
This is intentionally broader than the bad pure-remainder case. It can lose a
useful derivative guarantee when the absolute-value argument actually has a
proved eventual nonzero sign. It does not change the finite jet or its magnitude
bound. Exact output jets still receive exact-result treatment. The patch is a
correctness-first candidate, not a claim of a maximal regularity implementation.

Keeping this change in the ordinary \code{SeriesObservable} path also means that
replaying an observable recipe through that patched path applies the same
conservative policy. No external wrapper should attach a one-time warning while
leaving a replay route capable of silently restoring the old guarantee.

\subsection{The better implementation: track regularity inside the interpreter}
The more precise repair is to make the expression interpreter return both its
jet and an available smoothness order. For an absolute-value node, the transfer
rule is local and small:
\begin{equation}\label{eq:abs-policy}
 k_{\rm out}=\begin{cases}
 k_{\rm in},&\text{if the complete argument is eventually strictly positive or negative},\\
 \infty,&\text{if the complete argument is exactly zero},\\
 0,&\text{otherwise, using only the real Lipschitz estimate}.
 \end{cases}
\end{equation}
A nonempty leading block with a proved sign can provide the first case because
the stored error is asymptotically smaller than that block. An empty finite jet
cannot provide it merely by carrying a large derivative-order number.

The interpreter should aggregate this order through nested expressions, not
inspect only the outermost head. A magnitude-neutral intermediate \code{Abs}
can invalidate regularity before a later arithmetic operation. Conversely,
algebraic identities such as $|g|^2=g^2$ for real $g$ may restore smoothness, but
that is a separate proved identity, not a reason to retain the old order by
default. This narrowly scoped regularity transfer is the development opportunity
made concrete by N01.

\subsection{Acceptance cases and an optional weak-derivative extension}
The regression must inspect both the outer association and
\code{SeriesRepresentation}; a downgrade in only one place can be bypassed by
\code{seriesData}. The primary case is the public chain above. Controls should
include exact zero, an exact sign-known polynomial, a nonempty sign-known jet
with an inexact tail, nested absolute values, and replay of the observable recipe.
A sign-aware implementation should retain the useful order in the proved-sign
control; the conservative patch intentionally need not do so.

An optional future interface could distinguish classical derivatives from
almost-everywhere derivatives of real Lipschitz observables. That would require
an explicit change of contract. For this witness the first weak derivative is
locally bounded, but the distributional second derivative of $|g|$ also contains
point masses with weights $2|g'(y_n)|$. An ordinary power-log magnitude remainder
does not represent that discrete distribution. A weak-derivative mode therefore
must not be implemented by simply turning off the classical guard.

\section{N02: repeated affine probing makes a nonlinear fallback quadratic}
\subsection{A good local repair with a different global cost}
The new \code{certAffinePair} recognizes rational affine syntax by carrying a
pair $(a,b)$ for $ax+b$. Sums add their pairs. Products combine pairs while
refusing a product of two nonconstant affine factors. Other heads fail the
recognition attempt. One invocation visits each node of its admitted arithmetic
tree at most once, apart from exact rational arithmetic. The source comment's
linear-tree description is reasonable for that \emph{single invocation}.
\cite{cert}

The important integration point is different. Before the ordinary \code{Plus}
or \code{Times} interval fallback, \code{certEnclose} invokes
\code{certAffineRange}, which invokes \code{certAffinePair}. If that test fails,
\code{certEnclose} recursively processes the children, each of which may trigger
another full affine-recognition attempt. Successful affine roots are cheap; deeply
nested nonlinear roots repeatedly expose overlapping subtrees.
\begin{lstlisting}
MemberQ[{Plus, Times}, Head[expression]] &&
  (affine = certAffineRange[expression, x, interval]) =!= $Failed,
  certRoundInterval[affine, ctx],
Head[expression] === Plus,
  Fold[certAdd[#1, #2, ctx] &, {0, 0},
    certEnclose[#, x, interval, ctx] & /@ (List @@ expression)],
...
\end{lstlisting}
This integration is visible at \file{InverseCertificates.wl}, lines 103--148.
The finding concerns recognizer work before the old conservative fallback;
it does not question the soundness of an interval returned by that fallback.
\cite{cert}

\subsection{A canonical, supported witness family}
Use the unevaluated-by-expansion Horner family
\begin{equation}\label{eq:horner}
 h_0(x)=1+x^2,\qquad h_{d+1}(x)=1+xh_d(x).
\end{equation}
In Wolfram syntax this is
\begin{lstlisting}
expr = Nest[1 + x # &, 1 + x^2, depth];
\end{lstlisting}
The seed $1+x^2$ is deliberate. Starting from $x^2$ alone allows the first
multiplication to simplify to $x^3$, changing the base-case count. The chosen
seed leaves alternating \code{Plus} and \code{Times} nodes under ordinary
canonical evaluation without calling \code{Expand}. Every $h_d$ is a supported
polynomial, and the interval $[0,1/2]$ avoids any domain difficulty. There is no
unsupported special function at the bottom of the expression.

Let $A_d$ be the number of \code{certAffinePair} invocations in one recognition
attempt on $h_d$. The \code{Power} node $x^2$ is rejected immediately by the
affine recognizer, so
\[
 A_0=3,\qquad A_d=A_{d-1}+4=4d+3.
\]
For $d\ge1$, the intermediate product $x h_{d-1}$ costs
$1+1+A_{d-1}=4d+1$ recognizer invocations. The recognizer eagerly maps over the
children before it tests whether a failure occurred, so the failed nonlinear
subtree is traversed again at the next recursive interval level.

\begin{proposition}\label{prop:quadratic}
In the inspected recursive control structure, enclosing $h_d$ causes
\begin{equation}\label{eq:probe-count}
 C_d=4(d+1)^2-1
\end{equation}
invocations of the affine-pair recognizer. Thus recognizer work is quadratic in
$d$, even though the input syntax has linear size.
\end{proposition}
\begin{proof}
At depth zero, the only affine attempt is on $1+x^2$, so $C_0=3$.
For $d\ge1$, the evaluator first attempts the whole sum, then the product, then
recursively encloses $h_{d-1}$:
\[
 C_d=C_{d-1}+(4d+3)+(4d+1)=C_{d-1}+8d+4.
\]
Summing yields $C_d=3+4d^2+8d=4(d+1)^2-1$.
\end{proof}

This is a count of explicit recursive calls, not a claim that rational arithmetic
has constant bit cost. Integer/rational operand sizes, expression hashing,
list allocation, and interval arithmetic can add work. They do not eliminate
this repeated traversal in the source algorithm. The claim is also narrower
than a timing prediction for a complete certificate: a certificate evaluates
several expressions, including a derivative, and its surrounding proof work may
dominate a small example.

\subsection{Executed independent counts}
The supplied model constructs an explicit syntax tree with \code{Plus},
\code{Times}, an integer-power node, rational constants, and a variable. It
implements the inspected recognition/fallback control flow and compares it
with a single annotation pass. These are \emph{model operation counts}, not
native performance measurements.
\begin{center}
\begin{tabular}{@{}rrr@{}}
\toprule Horner depth $d$ & Repeated recognizer calls $C_d$ & Annotation nodes $4d+5$\\\midrule
4 & 99 & 21\\
8 & 323 & 37\\
16 & 1,155 & 69\\
32 & 4,355 & 133\\
64 & 16,899 & 261\\
128 & 66,563 & 517\\
\bottomrule
\end{tabular}
\end{center}
The annotation-node count uses the model's explicit literal exponent node; it
is not Wolfram's \code{LeafCount}. All tested depths matched the closed-form
recognizer count. The two evaluators also agreed on 300 randomized rational
arithmetic trees. Their interval outputs enclosed 2,100 exact rational sample
values. Sampling is a useful regression check, while interval soundness rests
on the elementary operation rules and their inductive composition.

\subsection{Repair with two passes, not more expansion}
A direct remedy is to annotate the expression tree once, then evaluate intervals
using the annotation. For each node retain either a rational affine pair or
\code{NotAffine}. All children are annotated bottom-up. Importantly, an
unsupported affine head such as \code{Power}, \code{Log}, or \code{Exp} does not
prevent useful affine annotations within its children.

The interval evaluator then behaves as follows. At a node carrying a pair, it
computes the affine endpoint range exactly and applies the existing outward
rounding at the same semantic point as today. At a node without a pair, it uses
the existing interval rule and the already annotated children. No polynomial
expansion, solving, or approximate coefficient test is required.

\begin{lemma}
For the rational arithmetic fragment implemented in the independent model,
replacing repeated recognition by bottom-up affine annotations preserves the
result of the original recognition-plus-interval algorithm.
\end{lemma}
\begin{proof}
Induct on the expression tree. The annotation computes exactly the same pair
recurrence as an independent recognition call; therefore both evaluators take
the affine shortcut at the same nodes and use the same exact endpoint formula.
At every other supported node, they apply the same interval operation to child
results that agree by induction.
\end{proof}

The production version should retain the existing rational/dyadic interval
primitives. The model deliberately does not reimplement the complete exponential
and logarithmic enclosure engine. Its purpose is to establish the structural
transformation and its local equivalence obligations, not to replace a trusted
certificate evaluator with an untested partial implementation.

A single pass returning both a pair and an interval is possible, but a careless
implementation computes rounded child intervals or raises child-domain failures
before discovering that the parent is exactly affine. The two-pass design cleanly
separates structural recognition from semantic interval evaluation and avoids
that ordering mistake. Annotation must remain arithmetic metadata, not a proof
that a nonaffine expression is continuous or in domain.

\subsection{Caching scope and implementation tradeoffs}
A cache can also remove repeated recognizer calls, but a cache keyed by complete
expression trees may still pay repeated structural hashing or equality costs.
Its actual cost depends on sharing and kernel representation. An explicit
annotated tree gives a straightforward linear structural-pass argument without
requiring such a representation assumption.

Affine annotations depend on the expression and chosen source symbol, but not
on the verification interval or enclosure precision. Interval values depend on
all of those. Mixing the two cache lifetimes would be incorrect: adaptive
refinement changes the interval and arithmetic context. A suitable implementation
can reuse a structural plan through several interval evaluations while recomputing
all interval-dependent values.

Annotation increases live structural storage by $\OO(N)$ for an $N$-node tree.
For tiny expressions, the existing local probe may be cheaper. A hybrid can keep
the simple path below a small structural threshold and use the annotated path
for larger trees, but the threshold should be determined by measurements rather
than guessed in this audit. The correctness invariant is the same in either
case: exact affine arithmetic remains before rounding.

\subsection{Native acceptance criteria}
The native runner supplies the stable Horner family at depths 8, 16, 32, and 64,
records actual \code{LeafCount}, and times \code{certEnclose} separately from
construction. It intentionally does not advertise an unmeasured speedup. A
production benchmark should additionally count recognizer invocations, with
instrumentation that distinguishes entry from exit tracing and does not itself
perform quadratic work.

Tests should compare old and planned enclosures on exactly affine, nested
nonlinear, mixed affine/transcendental, and unsupported expressions. Check the
same failures and interval inclusions, not only elapsed time. In particular,
retaining the affine repair's existing success cases is a regression obligation,
not a new finding in this article.

\section{N03: finding a discarded suffix by repeated membership}
\subsection{Where the new work occurs}
The recent \code{seriesTransportRemainderBound} routine obtains the source rows
and newly retained rows, then computes the discarded set as follows:
\begin{lstlisting}
before = d["Jet"][[1]];
after = r["SeriesRepresentation"]["Jet"][[1]];
removed = Select[before, ! MemberQ[after, #] &];
If[! SubsetQ[before, after], Return[result, Module]];
\end{lstlisting}
These are lines 516--518 of \file{SeriesOperations.wl}. The routine subsequently
uses the known discarded expression to transport an existing quantitative
bound. It also handles the case in which no retained block was removed.
\cite{operations}

The bound-transport theorem is already an earlier-review contribution and is
not being proposed again. The new issue is the selection algorithm used by the
implemented repair. Ordinary exponent truncation of a canonical ordered jet
normally retains an identical prefix. That common case does not require a
general membership query for every source row.

\subsection{What is established, and what is not}
Let the input have $N$ distinct rows and the result retain the first $M$ rows.
The source expression performs $N$ separate \code{MemberQ} calls on a list of
length $M$, followed by a subset check. \code{MemberQ} is a pattern-membership
operation; the official documentation does not promise an asymptotic complexity
or expose a reusable indexed set for these calls.\cite{wlmember}

Under a straightforward linear scan, the row comparisons in the first
statement alone are
\begin{equation}\label{eq:membership}
 \sum_{i=1}^{M}i+(N-M)M
 =NM-\frac{M(M-1)}2.
\end{equation}
For a no-op truncation, $M=N$, this is $N(N+1)/2$. This is the exact count in the
supplied linear-search model, not an instrumented count of the Wolfram kernel's
internal pattern matcher. An undocumented native optimization could change the
actual timing or comparison behavior. The source still issues repeated general
queries where a proved prefix relation is sufficient.

\begin{center}
\begin{tabular}{@{}rrr@{}}
\toprule Blocks, no-op case & Linear-membership model comparisons & Prefix row comparisons\\\midrule
16 & 136 & 16\\
64 & 2,080 & 64\\
256 & 32,896 & 256\\
512 & 131,328 & 512\\
\bottomrule
\end{tabular}
\end{center}
The table excludes the old subset check and the cost of comparing large
coefficient expressions. It therefore should not be read as a native wall-clock
speedup estimate. It identifies a local avoidable pattern and supplies a small
alternative whose structural work is explicit.

\subsection{A guarded prefix fast path}
The candidate replacement is:
\begin{lstlisting}
If[Length[after] <= Length[before] &&
    Take[before, Length[after]] === after,
  removed = Drop[before, Length[after]],
  removed = Select[before, ! MemberQ[after, #] &];
  If[! SubsetQ[before, after], Return[result, Module]]];
\end{lstlisting}
The length guard is important: it prevents an out-of-range \code{Take} from
becoming the new diagnostic behavior. The structural equality check is also
important. The patch does not merely assume that every stored representation
uses the same canonical coefficient syntax. If normalization has changed a row,
or the result is not a prefix, it falls back to the old selection and subset
policy.

\begin{proposition}\label{prop:prefix}
For a well-formed canonical jet with distinct rows, suppose
\[\texttt{after}=\operatorname{Take}(\texttt{before},M).\]
The proposed fast path returns the same discarded rows as the original membership expression, in the
same order, and the original subset test would succeed.
\end{proposition}
\begin{proof}
The first $M$ rows occur in \texttt{after}; every remaining distinct row does not.
The original \code{Select} therefore retains exactly the suffix beginning at
$M+1$. That is \code{Drop[before,M]}. Every row in the prefix belongs to the
source, so the subset relation follows without a second membership pass.
\end{proof}

This reasoning is deliberately restricted to canonical jets with distinct
rows. It is not a new validation rule for arbitrary user-built associations.
The fallback remains for other structurally admissible situations. It preserves
the original routine's handling rather than changing the meaning of membership
or attempting to use a symbolic inequality between coefficient expressions.

\subsection{Why this is a useful low-risk patch}
The fast path affects only the extraction of the known discarded finite part.
The subsequent bound formula, conditions, signed-bound policy, contract record,
and presentation of integer Dirichlet powers remain untouched. This keeps the
change separable from mathematical bound improvements already in the backlog.

A semantic no-op is a particularly useful benchmark because it isolates a case
where no discarded contribution is needed. The candidate still compares a
prefix, but avoids re-searching it for each row. For a genuine truncation it
extracts the suffix directly. Its cost is linear in the structural size of the
rows inspected or copied; coefficient-expression bit/term complexity is not
assumed constant.

The independent suite checked 500 randomly generated prefix cases, reordered
and nonprefix fallbacks, empty lists, complete retention, and refusal of an
invalid subset in the model. The native runner additionally specifies no-op
truncations of public Zeta expansions. Those native observations have not been
executed here. Their actual constructor and truncation costs may include other
work, so the helper and the whole operation should be profiled separately.

\section{Implementation package and focused development sequence}
\subsection{What the supplied artifacts implement}
The archive contains a standard-library Python model and an executable test
runner. The affine model includes both repeated recognition and a bottom-up
annotation plan. The truncation model includes the repeated linear-membership
algorithm and the guarded prefix alternative. The regularity witness is checked
through an exact derivative-coefficient recurrence, elementary rational bounds,
and supplementary numerical one-sided difference quotients.

The patch emitter produces two candidate edits to the canonical
\file{src/Kernel/SeriesOperations.wl}: the conservative \code{Abs} regularity
downgrade, and the guarded prefix fast path. It does not modify the repository.
It verifies the audited Git commit by default and requires unique exact source
anchors. An explicit override permits an unversioned source tree while retaining
anchor validation. Its fixture tests cover both edits independently, missing and
duplicate anchors, already changed input, and preservation of unrelated text.

No production \code{InverseCertificates.wl} patch is supplied. N02 has an
implemented independent reference design and a native benchmark specification,
but integrating the annotation plan into all certificate expression rules and
failure paths deserves a full native validation pass. Supplying a partial
certificate-engine replacement as if it were complete would be misleading.

\subsection{Ordering the work}
\begin{longtable}{@{}L{12mm}L{62mm}L{\dimexpr\textwidth-74mm-4\tabcolsep\relax}@{}}
\toprule Step & Change & Completion criterion\\\midrule\endhead
1 & Apply a conservative regularity downgrade for inexact absolute-value observables. & The N01 chain no longer obtains a classical derivative from the transported input contract; exact controls remain correct.\\
2 & Refine that repair into sign-aware regularity transfer inside the expression interpreter. & Proved-sign cases preserve their legitimate order; unknown zero-crossing cases do not; nested observables and recipe replay agree.\\
3 & Install the guarded prefix selection path. & Prefix/no-op cases select the same discarded rows and preserve the existing quantitative-bound semantics; nonprefix cases retain the fallback.\\
4 & Introduce a certificate structural annotation plan. & Affine decisions and outward interval behavior match the old evaluator while Horner-family recognition work becomes linear in tree size.\\
5 & Add narrow scaling and semantic-composition regressions. & Structural work counts and regularity transitions are tested separately from numerical values and total timings.\\
\bottomrule
\end{longtable}

This sequence is intentionally local. It does not reproduce the already extensive
roadmap for additional scales, special functions, native-backend coverage,
integration, general proof objects, or multivariate systems. The opportunities
identified here are operation-specific regularity transfer, reusable
interval-independent affine plans, and exploiting the exact prefix invariant
of truncation.

\subsection{Tests should target the interaction, not only the repaired feature}
N01 is invisible to tests that inspect only \code{Normal[result]} and its
remainder power. The invalid fact is in the derivative capability consumed by a
later operation. The essential regression is therefore a sequence:
construction with an honest remainder family, arithmetic cancellation,
differentiation with a legitimate declaration, absolute value, and a second
classical differentiation.

N02 is invisible to tests containing only successfully affine expressions.
A recognizer that is fast when it succeeds can be expensive when it repeatedly
fails just above an old nonlinear subtree. The depth family deliberately holds
the mathematical fragment fixed and changes only nesting.

N03 is invisible to tests of bound values alone. The implemented bound can be
perfectly valid while its no-op transport makes unnecessary general searches.
The acceptance suite should retain both the semantic checks and a narrowly
scoped structural-work observation. A total-time threshold alone is too noisy
to explain which phase regressed.

These three test shapes are distinct from a blanket request for more tests.
They are specific interaction witnesses and size families derived from the
current implementation.

\section{Reproduction and recorded results}
\subsection{Independent checks that were executed}
Run from the archive directory:
\begin{lstlisting}
python code/verify_independent.py
python code/test_patch_emitter.py
\end{lstlisting}
Python 3.10 or later and its standard library are sufficient. The first command
writes \file{evidence/independent_results.json} and its text transcript. The
second writes the patch-fixture result record and transcript. The delivered
records come from actual executions during this audit.

\begin{tabularx}{\textwidth}{@{}lXr@{}}
\toprule Suite & What was checked & Result\\\midrule
Independent model tests & Derivative recurrence, cusp witnesses, inverse-family bounds, affine recognition, exact interval equivalence, prefix/fallback selection, and structural count formulas. & 14/14\\
Patch-emitter fixtures & Unique-anchor transformations, individual edits, missing/duplicate anchors, and refusal to reapply edits. & 6/6\\
\bottomrule
\end{tabularx}

The independent suite includes 300 randomized expression trees, 2,100 exact
rational evaluation points inside their intervals, and 500 randomized prefix
cases. These are subcases within the 14 tests, not thousands of package test
passes. Random seeds are fixed in the source. Mathematical proofs in the article
supply the general statements; the finite samples are supplementary executable
checks.

The cusp samples numerically approximate left and right difference quotients
normalized by $y_n$. They approach $-\sqrt{10}$ and $+\sqrt{10}$. This numerical
calculation is not used to infer non-differentiability: the exact local expansion
in Proposition~\ref{prop:cusp} already proves it.

\subsection{Native observations that remain to be run}
Use a fresh kernel and an existing checkout of the pinned revision. For example,
in PowerShell:
\begin{lstlisting}
$env:ASYMPTOTIC_REPO = 'C:\src\Asymptotic'
wolframscript -file code/run_native.wls
\end{lstlisting}
For a POSIX shell:
\begin{lstlisting}
ASYMPTOTIC_REPO=/path/to/Asymptotic \
  wolframscript -file code/run_native.wls
\end{lstlisting}
The script reads the canonical modular loader. It records the kernel version,
the observed Git commit when available, the N01 public chain, an exact
sign-known control, Horner-family interval timings, and Zeta no-op truncation
observations. Individual observations have a bounded execution time. It does
not download code, install dependencies, edit the repository, or claim to run
the upstream complete suite.

The script's exit status distinguishes a runner/load/output failure from
successfully writing observations. A zero exit code is \emph{not} a correctness
verdict: an observation can show the baseline defect, a timeout, a message, or a
rejected operation. Inspect the serialized results. The delivered archive does
not contain a fabricated \file{native_observations.json}; that file is created
only by an actual native execution.

The runner is a Wolfram-kernel observation tool, not a promise of compatibility
with every Mathics release. The mathematical counterexample and the structural
algorithms do not depend on a particular kernel, but public routing, native
\code{Series}, pattern matching, and execution timing do. A Mathics validation
should record its own results separately rather than relabeling Wolfram or
Python results.

\subsection{Emitting and reviewing the candidate patch}
\begin{lstlisting}
python patches/emit_candidate_patch.py /path/to/Asymptotic \
  > candidate.diff
\end{lstlisting}
Use \code{--which regularity} or \code{--which prefix} to emit one change only.
The emitter is read-only. It refuses an unexpected Git commit by default and
refuses missing or duplicate anchors. The explicit
\code{--skip-commit-check} option is for a consciously selected unversioned or
modified source tree; it does not disable the anchor checks.

The source anchors were compared with the pinned reads. The transformation was
fixture-tested, but it was not applied to a complete local upstream checkout in
this audit. The generated patch still needs review, native execution of the
focused cases, and the repository's normal validation. Rebuild the standalone
using \file{validation/build_standalone.py}; editing the generated distribution
alone would not update the canonical source.\cite{layout}

The regularity patch is intentionally conservative and can reduce supported
differentiation in some valid sign-known inexact cases. The prefix patch has a
narrower semantic footprint and retains the old fallback. Neither is presented
as a completed production release.

\section{Conclusion}
The strongest new result is N01. The existing derivative guard can be obeyed
honestly at the input and still be bypassed indirectly by an operation that
preserves magnitude but not differentiability. The example is not based on a
complex branch, a malformed object, or a missing initial derivative estimate.
Its infinitely many simple zeros explain exactly which premise is lost.

N02 and N03 concern the implementation of two useful recent fixes rather than
a rejection of those fixes. Affine cancellation should remain exact before
rounding, but the recognition plan need not be rediscovered at every nonlinear
subtree. Quantitative truncation bounds should remain preserved, but a known
prefix can be removed without repeated general membership searches.

The immediate actionable changes are small: downgrade unproved regularity in
the absolute-value observable path and add the guarded prefix fast path. The
larger follow-on is a reusable, interval-independent affine annotation plan.
Together they give specific correctness and performance improvements without
repeating the repository's established development backlog.

\appendix
\section{Source coverage and exclusions}\label{app:coverage}
The table is a source-reading record, not a test-coverage percentage. ``Selected''
means substantial relevant windows rather than a verified complete-file read.
Several large connector responses were windowed or truncated, so completeness
is not inferred from a request for an entire file.

\begin{longtable}{@{}L{65mm}L{20mm}L{\dimexpr\textwidth-85mm-4\tabcolsep\relax}@{}}
\toprule Source under \file{src/Kernel/} & Read scope & Focus and disposition\\\midrule\endhead
\file{AsymptoticAnalysis.wl} & Selected & Public inverse assembly, precision-tracked jets, arithmetic, \code{fwdAbs}, representation, formatting, and residual dispatch. N01 uses the absolute-value branch.\\
\file{SeriesOperations.wl} & Selected & Representation normalization, observables, composition, truncation, derivative guard, and refinement entry paths. N01 and N03 retained.\\
\file{InverseCertificates.wl} & Selected & Rational interval primitives, affine recognition, recursive enclosure, source conditions, seed and adaptive certificate flow. N02 retained.\\
\file{SeriesEnvelopeArithmetic.wl} & Substantial & Composite-envelope arithmetic and approach handling. No additional finding retained.\\
\file{SourceCoordinates.wl} & Selected & Logarithmic/exponential source charts, observable restoration, and error transport. No additional finding retained.\\
\file{GammaForward.wl} & Substantial & Positive-factor logarithmic normalization and forward carrier extraction. No additional finding retained.\\
\file{GammaInverse.wl} & Substantial & Ordered coefficient recurrence, affine parameters, omitted block bound, and truncation. No additional finding retained.\\
\file{SpecialFunctionAdapters.wl} & Selected & Erfc, Gamma/LogGamma, and local threshold model construction. No additional finding retained.\\
\file{ParameterizedSpecialFunctions.wl} & Substantial & Fixed-parameter analytic and Frobenius adapters. No additional finding retained.\\
\file{DirichletSpecialFunctions.wl} & Substantial & Zeta/Lerch block construction and explicit bound metadata; supports the N03 benchmark specification.\\
\file{NumericalInverseChecks.wl} & Substantial & Stored-equation scope and numerical source-domain checks. No additional finding retained.\\
\file{ExactTermination.wl} & Substantial & Bounded symbolic verification of inverse termination. No additional finding retained.\\
\file{NativeCompatibility.wl} & Selected & Native-result distinction and backend/held request entry; checked the reproduction option spelling. No additional finding retained.\\
\file{MathicsAlgebra.wl} & Substantial & Algebraic substitutes and list helpers; previously catalogued concerns excluded.\\
\file{MathicsAssumptions.wl} & Substantial & Realness/sign proof helpers and assumption walking. No additional finding retained.\\
\file{MathicsCalculus.wl} & Substantial & Perturbative inverse expression and logarithmic Euler derivative. No additional finding retained.\\
\file{MathicsSimplification.wl} & Substantial & Simplification, assumptions, and Series/Limit wrappers. Previously catalogued concerns excluded.\\
\file{MathicsTaylor.wl} & Substantial & Restricted defining Taylor series. Previously catalogued parameter restrictions excluded.\\
\file{MathicsSpecialFunctions.wl} & Substantial & Finite Stirling forward adapter. No additional finding retained.\\
\bottomrule
\end{longtable}

The ordinary calculus tests in \file{src/Tests/SeriesOperations.wlt} were also
read, notably their explicit derivative-contract tests and exact finite-tail
control. They were not executed. Architecture descriptions of other modules
were read in the kernel README, but that is not counted as inspection of those
modules' implementation.\cite{optests,layout}

The maintained review register was used to exclude already catalogued work even
when a related code path remained interesting. Consequently, the absence of a
new finding in the table should not be interpreted as a claim that every
existing concern is resolved. That would contradict the purpose of a
nonduplicating incremental audit.

\section{Artifact inventory}
\begin{longtable}{@{}L{71mm}L{\dimexpr\textwidth-71mm-2\tabcolsep\relax}@{}}
\toprule File & Role\\\midrule\endhead
\file{article.tex}, \file{article.pdf} & Editable article and rendered version.\\
\file{README.md} & Reproduction commands and execution boundaries.\\
\file{code/audit_models.py} & Independent rational/structural reference models and derivative recurrence.\\
\file{code/verify_independent.py} & Executed 14-test suite; writes its actual result records.\\
\file{code/test_patch_emitter.py} & Executed six-test synthetic fixture suite.\\
\file{code/run_native.wls} & Unexecuted native observations; records actual results when run locally.\\
\file{code/check_wl_lexical.py} & Delimiter, string, and nested-comment balance check; not a native parser.\\
\file{patches/emit_candidate_patch.py} & Read-only emitter for the two conservative candidate edits.\\
\file{evidence/independent_results.json} & Actual independent counts, subcase totals, and cusp samples.\\
\file{evidence/independent_tests.txt} & Actual independent-test transcript.\\
\file{evidence/patch_fixture_results.json} & Actual patch-fixture outcome.\\
\file{evidence/patch_fixture_tests.txt} & Actual fixture-test transcript.\\
\file{evidence/audit_manifest.json} & Snapshot, scope, finding status, and execution record.\\
\file{evidence/novelty_crosswalk.csv} & Concise mapping to nearest previous entries without reissuing them.\\
\file{evidence/source_inventory.csv} & Pinned source-reading inventory.\\
\file{evidence/wl_lexical_check.json} & Actual lexical-balance check outcomes.\\
\file{build.sh} & PDF build command sequence.\\
\bottomrule
\end{longtable}
No upstream package distribution, font files, checksum files, or fabricated
native result files are included.

\begin{thebibliography}{99}
\small
\bibitem{repo} Vladimir Reshetnikov. \emph{Asymptotic}, audited commit
\code{651f2029d0b2cd4da9e4dfdf1f4275a124d23b99}.
\url{https://github.com/VladimirReshetnikov/Asymptotic/tree/651f2029d0b2cd4da9e4dfdf1f4275a124d23b99}.
Repository README and commit metadata, accessed 10 September 2026.

\bibitem{layout} \emph{Canonical kernel layout}.
\url{https://github.com/VladimirReshetnikov/Asymptotic/blob/651f2029d0b2cd4da9e4dfdf1f4275a124d23b99/src/Kernel/README.md}.

\bibitem{reviewindex} \emph{Imported code-review index}.
\url{https://github.com/VladimirReshetnikov/Asymptotic/blob/651f2029d0b2cd4da9e4dfdf1f4275a124d23b99/external-reports/code-review/README.md}.

\bibitem{status} \emph{Code review implementation status}.
\url{https://github.com/VladimirReshetnikov/Asymptotic/blob/651f2029d0b2cd4da9e4dfdf1f4275a124d23b99/docs/development/CODE_REVIEW_STATUS.md}.

\bibitem{wave3} \emph{Wave 3 intake}.
\url{https://github.com/VladimirReshetnikov/Asymptotic/blob/651f2029d0b2cd4da9e4dfdf1f4275a124d23b99/docs/development/WAVE_3_INTAKE.md}.

\bibitem{wave4} \emph{Wave 4 intake}.
\url{https://github.com/VladimirReshetnikov/Asymptotic/blob/651f2029d0b2cd4da9e4dfdf1f4275a124d23b99/docs/development/WAVE_4_INTAKE.md}.

\bibitem{review14} \emph{Incremental technical audit}, imported review 14.
Affine and truncation sections, especially source lines 263--355 and 470--580.
\url{https://github.com/VladimirReshetnikov/Asymptotic/blob/651f2029d0b2cd4da9e4dfdf1f4275a124d23b99/external-reports/code-review/wave-2/code-review-14/article.tex}.

\bibitem{affinecommit} \emph{Evaluate rational affine subexpressions exactly in certificates (C20)},
commit of 10 September 2026, 13:37:49 UTC.
\url{https://github.com/VladimirReshetnikov/Asymptotic/commit/e9142f7f5e3db241d5ffb8951c72606ed0c89aa9}.

\bibitem{kernel} \emph{AsymptoticAnalysis.wl}, canonical main module.
Relevant definitions: \code{fwdAbs}, precision-tracked jets, inverse assembly,
and \code{GeneralizedSeries} accessors.
\url{https://github.com/VladimirReshetnikov/Asymptotic/blob/651f2029d0b2cd4da9e4dfdf1f4275a124d23b99/src/Kernel/AsymptoticAnalysis.wl}.

\bibitem{operations} \emph{SeriesOperations.wl}.
Relevant source windows: 305--374 (observable interpretation and construction),
506--558 (bound transport and differentiation), and the representation helpers
near the beginning of the module.
\url{https://github.com/VladimirReshetnikov/Asymptotic/blob/651f2029d0b2cd4da9e4dfdf1f4275a124d23b99/src/Kernel/SeriesOperations.wl}.

\bibitem{cert} \emph{InverseCertificates.wl}.
Relevant source windows: 95--151 (affine recognizer and enclosure dispatch),
251--327 (certificate attempt), and 357--484 (public certificate orchestration).
\url{https://github.com/VladimirReshetnikov/Asymptotic/blob/651f2029d0b2cd4da9e4dfdf1f4275a124d23b99/src/Kernel/InverseCertificates.wl}.

\bibitem{numerical} \emph{NumericalInverseChecks.wl}.
\url{https://github.com/VladimirReshetnikov/Asymptotic/blob/651f2029d0b2cd4da9e4dfdf1f4275a124d23b99/src/Kernel/NumericalInverseChecks.wl}.

\bibitem{optests} \emph{SeriesOperations.wlt}, upstream ordinary calculus tests.
\url{https://github.com/VladimirReshetnikov/Asymptotic/blob/651f2029d0b2cd4da9e4dfdf1f4275a124d23b99/src/Tests/SeriesOperations.wlt}.

\bibitem{wlabs} Wolfram Research. \emph{Abs}, Wolfram Language documentation.
\url{https://reference.wolfram.com/language/ref/Abs.html}.
Accessed 10 September 2026. Only the real absolute-value meaning is used here.

\bibitem{wlmember} Wolfram Research. \emph{MemberQ}, Wolfram Language documentation.
\url{https://reference.wolfram.com/language/ref/MemberQ.html}.
Accessed 10 September 2026. Documents pattern-membership semantics, not a
complexity guarantee for the audit's repeated-list-query pattern.
\end{thebibliography}
\end{document}
